\section{Introduction}\label{sec:introduction}

Accurate prediction of gene expression from DNA sequence and chromatin context is a longstanding goal in computational genomics. In crops, where experimental profiling across tissues, genotypes, and conditions remains costly, computational models that generalize across biological contexts are particularly valuable. Wheat (\textit{Triticum aestivum}), with its large hexaploid genome (16\,Gb) and 117,319 annotated genes, poses special challenges for expression modelling.

Existing approaches such as DeepWheat \citep{DeepWheat2024} train separate convolutional neural networks for each tissue, each requiring 24--54M parameters. While effective for leaf (PCC\,=\,0.864) and spike (PCC\,=\,0.820), performance degrades substantially for developmental tissues when evaluated on an independent held-out test split (mean PCC across six tissues\,=\,0.461), and the approach cannot share learned representations across tissues.

Genomic foundation models---large pretrained models that learn general-purpose sequence representations from unlabelled genomes---offer a promising alternative. Models such as Nucleotide Transformer \citep{DallaTorre2023}, DNABERT-2 \citep{Zhou2024}, and the plant-specific \botanic \citep{Botanic2025} have shown strong performance on variant effect prediction and regulatory element classification. However, their application to multi-tissue expression prediction has been limited.

Here we present \mytool, a unified multi-tissue expression predictor that leverages LoRA-adapted \citep{Hu2022} embeddings from \botanic (318M parameters, BiMamba2 architecture) combined with epigenomic mark features via a lightweight prediction head with FiLM conditioning \citep{Perez2018}. \mytool predicts expression for all six wheat tissues with a single 1.0M-parameter head, achieving PCC\,=\,0.910 on log-transformed expression---a 22\% improvement over the best frozen-embedding baseline and a 97\% improvement over the mean DeepWheat performance.
